\section*{Capítulo \ref{ch:g7:atoms-and-molecules} --- Átomos y moléculas}

\begin{solution}{exo:g7:atoms-and-molecules:1}
Un \omterm{def:g7:atoms-and-molecules:atom}{átomo} es una de las partículas extremadamente pequeñas con las que
está construida toda la materia (existen un centenar de clases). Una
\omterm{def:g7:atoms-and-molecules:molecule}{molécula} es un grupo de \omterm{def:g7:atoms-and-molecules:atom}{átomos} fuertemente unidos, el trozo más pequeño
de una sustancia que sigue siendo esa sustancia.
\end{solution}

\begin{solution}{exo:g7:atoms-and-molecules:2}
\ce{H}, \ce{C}, \ce{N}, \ce{O}, \ce{Na}, \ce{Fe}.
\end{solution}

\begin{solution}{exo:g7:atoms-and-molecules:3}
1 \omterm{def:g7:atoms-and-molecules:atom}{átomo} de carbono y 2 \omterm{def:g7:atoms-and-molecules:atom}{átomos} de oxígeno: 3 \omterm{def:g7:atoms-and-molecules:atom}{átomos} en total.
\end{solution}

\begin{solution}{exo:g7:atoms-and-molecules:4}
Agua, oxígeno, metano, nitrógeno.
\end{solution}

\begin{solution}{exo:g7:atoms-and-molecules:5}
\ce{NO2}.
\end{solution}

\begin{solution}{exo:g7:atoms-and-molecules:6}
\ce{O2} es una \omterm{def:g7:atoms-and-molecules:molecule}{molécula} formada por dos \omterm{def:g7:atoms-and-molecules:atom}{átomos} de oxígeno unidos;
\ce{2O} son dos \omterm{def:g7:atoms-and-molecules:atom}{átomos} de oxígeno separados. \ce{CO} es una \omterm{def:g7:atoms-and-molecules:molecule}{molécula} de
un \omterm{def:g7:atoms-and-molecules:atom}{átomo} de carbono y uno de oxígeno (dos mayúsculas, dos \omterm{def:g7:atoms-and-molecules:symbol}{símbolos});
\ce{Co} es un \omterm{def:g7:atoms-and-molecules:atom}{átomo} de cobalto (un solo \omterm{def:g7:atoms-and-molecules:symbol}{símbolo}, mayúscula y luego
minúscula).
\end{solution}

\begin{solution}{exo:g7:atoms-and-molecules:7}
\ce{3H2O}: $3\times 2 = 6$ \omterm{def:g7:atoms-and-molecules:atom}{átomos} de hidrógeno y $3\times 1 = 3$ \omterm{def:g7:atoms-and-molecules:atom}{átomos}
de oxígeno. \ce{5O2}: ningún hidrógeno, $5\times 2 = 10$ \omterm{def:g7:atoms-and-molecules:atom}{átomos} de
oxígeno.
\end{solution}

\begin{solution}{exo:g7:atoms-and-molecules:8}
El negro es el carbono y el rojo, el oxígeno: un carbono, dos oxígenos,
\ce{CO2}, el dióxido de carbono.
\end{solution}

\begin{solution}{exo:g7:atoms-and-molecules:9}
El agua de Dalton era \ce{HO}. La \omterm{def:g7:atoms-and-molecules:molecule}{molécula} real, \ce{H2O}, tiene un
\omterm{def:g7:atoms-and-molecules:atom}{átomo} de hidrógeno más: a la \omterm{def:g7:atoms-and-molecules:molecule}{molécula} de Dalton le faltaba un \omterm{def:g7:atoms-and-molecules:atom}{átomo}.
\end{solution}

\begin{solution}{exo:g7:atoms-and-molecules:10}
El agua pura no es una \omterm{def:g6:pure-substances-and-mixtures:mixture}{mezcla}: todas sus \omterm{def:g7:atoms-and-molecules:molecule}{moléculas} son \omterm{def:g7:atoms-and-molecules:molecule}{moléculas}
idénticas de \ce{H2O}. El \omterm{def:g4:air-a-mixture-of-gases:air}{aire} es una \omterm{def:g6:pure-substances-and-mixtures:mixture}{mezcla}: contiene \omterm{def:g7:atoms-and-molecules:molecule}{moléculas} de
nitrógeno, \omterm{def:g7:atoms-and-molecules:molecule}{moléculas} de oxígeno, \omterm{def:g7:atoms-and-molecules:atom}{átomos} de argón y unas pocas \omterm{def:g7:atoms-and-molecules:molecule}{moléculas}
de dióxido de carbono.
\end{solution}

\begin{solution}{exo:g7:atoms-and-molecules:11}
Diez millones de \omterm{def:g7:atoms-and-molecules:atom}{átomos} forman \qty{1}{mm}, así que
\qty{1}{cm} $= \qty{10}{mm}$ requiere $10 \times$ diez millones $=$
cien millones de \omterm{def:g7:atoms-and-molecules:atom}{átomos}. $\qty{100}{m} = \num{100000}$~mm requiere $\num{100000}\times$
diez millones $=$ un billón de \omterm{def:g7:atoms-and-molecules:atom}{átomos} (un millón de millones).
\end{solution}

\begin{solution}{exo:g7:atoms-and-molecules:12}
$6 + 12 + 6 = 24$ \omterm{def:g7:atoms-and-molecules:atom}{átomos}. Hidrógeno: $\frac{12}{24} = \qty{50}{\percent}$.
Carbono: $\frac{6}{24} = \qty{25}{\percent}$ (y el oxígeno, el otro
\qty{25}{\percent}).
\end{solution}

\begin{solution}{pb:g7:atoms-and-molecules:1}
\textbf{1.} Nitrógeno \ce{N2}, 2 \omterm{def:g7:atoms-and-molecules:atom}{átomos}; oxígeno \ce{O2}, 2 \omterm{def:g7:atoms-and-molecules:atom}{átomos};
dióxido de carbono \ce{CO2}, 3 \omterm{def:g7:atoms-and-molecules:atom}{átomos}.

\textbf{2.} Los \omterm{def:g7:atoms-and-molecules:atom}{átomos} de argón no se unen formando \omterm{def:g7:atoms-and-molecules:molecule}{moléculas}: el argón
está formado por \omterm{def:g7:atoms-and-molecules:atom}{átomos} sueltos.

\textbf{3.} Nitrógeno: $78 \times 2 = 156$ \omterm{def:g7:atoms-and-molecules:atom}{átomos}. Oxígeno: $21\times 2 =
42$ \omterm{def:g7:atoms-and-molecules:atom}{átomos}.

\textbf{4.} $156 + 42 + 1 = 199$ \omterm{def:g7:atoms-and-molecules:atom}{átomos}.

\textbf{5.} $\frac{156}{199} \approx 0.78$, alrededor del
\qty{78}{\percent}.

\textbf{6.} En \omterm{def:g7:atoms-and-molecules:molecule}{moléculas} de oxígeno: $16\times 2 = 32$; en \omterm{def:g7:atoms-and-molecules:molecule}{moléculas} de
dióxido de carbono: $4 \times 2 = 8$; en total, $32 + 8 = 40$.

\textbf{7.} Faltan $42 - 40 = 2$ \omterm{def:g7:atoms-and-molecules:atom}{átomos} de oxígeno, el equivalente a una
\omterm{def:g7:atoms-and-molecules:molecule}{molécula} de oxígeno. Han acabado en \omterm{def:g7:atoms-and-molecules:molecule}{moléculas} de agua, \ce{H2O}, que el
cuerpo fabrica con el hidrógeno de los alimentos; esa agua sale en parte
como vapor de agua en el aliento (eliminado en estas cifras).

\textbf{8.} En el \omterm{def:g4:air-a-mixture-of-gases:air}{aire} inspirado: 0. En el espirado: 4 (uno en cada
\ce{CO2}).

\textbf{9.} De los alimentos que usa el cuerpo: los azúcares y las
grasas contienen \omterm{def:g7:atoms-and-molecules:atom}{átomos} de carbono, y el cuerpo los libera en el
dióxido de carbono que espira.

\textbf{10.} Han desaparecido $21 - 16 = 5$ \omterm{def:g7:atoms-and-molecules:molecule}{moléculas} de oxígeno y han
aparecido 4 \omterm{def:g7:atoms-and-molecules:molecule}{moléculas} de dióxido de carbono.

\textbf{11.} Nitrógeno: dos bolas azules unidas por tres varillas.
Oxígeno: dos bolas rojas unidas por dos varillas. Dióxido de carbono:
una bola negra entre dos bolas rojas, cada una unida a ella por dos
varillas, en línea recta. Argón: una sola bola azul claro.

\textbf{12.} $156 + 32 + 12 + 1 = 201$ \omterm{def:g7:atoms-and-molecules:atom}{átomos} (las 4 \omterm{def:g7:atoms-and-molecules:molecule}{moléculas} de
dióxido de carbono contienen 12 \omterm{def:g7:atoms-and-molecules:atom}{átomos}). Son $201 - 199 = 2$ \omterm{def:g7:atoms-and-molecules:atom}{átomos} más
que en el \omterm{def:g4:air-a-mixture-of-gases:air}{aire} inspirado: 4 \omterm{def:g7:atoms-and-molecules:atom}{átomos} de carbono ganados de los alimentos,
menos los 2 \omterm{def:g7:atoms-and-molecules:atom}{átomos} de oxígeno que se fueron en el agua.
\end{solution}
